\scanpage{227}
\mspage{177}
% L1
Donc il faut prouver simplement ceci

% L2–5
\uline{Lemme} Supposons que $\pi_i(f,x)$ iso pour $i\in[0,n-1]$, et b) satisfait pour $i\in[0,n-1]$, où $n\in\mathbb N$, $n\geq\unsure{2|1}$, est fixé. Alors les conditions suivantes sont équivalentes :
% L5 margin
% slanted
\margin{Sans doute faux}
% L6
a) $\pi_n(\deleted{\unsure{x}}f,x)$ \inserted{$:\pi_n(X,x)\unsure{-|\to}\pi_n(Y,\unsure{(u)|fx|(x)})$} iso pour tout $x\in\mathrm{Ob}\,X$

% L7–8
b) $H^n(Y,F)\xrightarrow{\sim}H^n(X,f^*F)$ iso pour tt faisceau abélien loc. constant sur $Y$.

% L9–10
Notons que \deleted{$\pi_{\unsure{i}}(f,x)=\pi_{i-1}($} $f$ induit une flèche \unsure{en|au} $\mathrm{Cat}_{\bullet}$
% L11
\[
\Omega(f,x):\Omega(X,x)\longrightarrow\Omega(Y,fx)
\]
% L12
d'où
% L12–14
\diagram{\begin{tikzcd}[ampersand replacement=\&,row sep=small]
\pi_i(\Omega(X,x).) \arrow[r] \arrow[d,phantom,"{\simeq}"{description,rotate=90}] \& \pi_i(\Omega(Y,fx) \arrow[d,phantom,"{\simeq}"{description,rotate=90}]\\
\pi_{i+1}(X,x)\arrow[r] \& \pi_{i+1}(Y,fx
\end{tikzcd},}
% L15
s'identifiant à $\pi_{i+1}(f,x)$ :
% L16
\[
\pi_i(\Omega(f,x))\simeq\pi_{i+1}(f,x).
\]
% L17
L'hypothèse \deleted{implique que} a) \inserted{du lemme} implique que
% L18
\[
\pi_i(\Omega(\deleted{\unsure{f|X}}\unsure{X|f},x))\xrightarrow{\sim}\pi_i(\Omega(Y,y))\qquad\text{pour }i\in[0,n-1].
\]
% L19–25
Il s'agit ici des $\pi_i$ basés en les points \unsure{unités}. On devra sans doute prouver que, du fait qu'il s'agit d'un morphisme de monoïdes dans $\mathrm{Cat}$ (satisfaisant certaines conditions ?) ça reste vrai pour les $\pi_i$ en \uline{tout} point $\xi$ de $\Omega(X,x)$. Par l'hyp. de récurrence, on saura donc que $\Omega(X,x)\to\Omega(Y,y)$
% archive sheet number: 227
